\documentclass[11pt]{article}
\usepackage[T1]{fontenc}
\usepackage{lmodern}
\usepackage[margin=0.88in]{geometry}
\usepackage{amsmath,amssymb,amsthm}
\usepackage{booktabs,longtable,array}
\usepackage[hidelinks]{hyperref}
\hypersetup{
  pdftitle={An Exact Observation Interface for the Normobaric Oxygen Paradox},
  pdfauthor={Anthony Vito Coppa},
  pdfsubject={Metric Mirror Coordinates, Adaptation Memory, and Empirical Traceability}
}
\setlength{\parindent}{0pt}
\setlength{\parskip}{4pt}
\setlength{\emergencystretch}{2em}
\renewcommand{\arraystretch}{1.12}
\newtheorem{theorem}{Theorem}
\newtheorem{proposition}[theorem]{Proposition}
\newtheorem{corollary}[theorem]{Corollary}
\newcommand{\R}{\mathbb R}
\newcommand{\pos}{\mathbb R_{>0}}
\newcommand{\Id}{I}
\newcommand{\Cr}{C_{\mathrm{redox}}}
\newcommand{\code}[1]{\nolinkurl{#1}}
\newcommand{\OPEN}{\textup{OPEN}}
\title{\Large\bfseries An Exact Observation Interface\\
for the Normobaric Oxygen Paradox\\[7pt]
\normalsize\mdseries Metric Mirror Coordinates, Adaptation Memory,\\
and Empirical Traceability}
\author{Anthony Vito Coppa\\\small Coppalini Architecture}
\date{Version 0.2.0 \quad 12 November 2026}

\begin{document}
\maketitle
\vspace{-1em}
\begin{abstract}
We construct an invertible interface from a declared pair of positive oxidant
and antioxidant assay readings to the established Metric Mirror operators.
With the paired baselines retained, contrast and magnitude recover both
original numerical readings exactly. An explicit matrix intertwines exchange
of the two normalized coordinates with the root mirror and converts their
logarithmic contrast into the axial parameter. A sampled adaptation equation
gives an exact return-with-memory theorem, including the input jump,
interval composition, and finite recovery. Positive interval arithmetic carries
the entered uncertainty through the same interface. The empirical connection
is specified field by field for primary normobaric oxygen paradox studies.
The completed result is a reproducible observation and model interface:
its coordinate identities and adaptation theorem are proved here, while
participant-level parameter identification and the ROS--HIF causal law remain
\OPEN. All numerical demonstrations supplied with the workbench are synthetic.
\end{abstract}

\section{The observation and its owning reference}

The normobaric oxygen paradox (NOP) literature studies responses after a
hyperoxic exposure returns to a lower oxygen condition. The proposed
ROS/scavenger mechanism invokes different recovery times; the linked review
also records conflicting roles of ROS in HIF regulation \cite{review}.
The present construction asks a directly testable question:
\emph{when the declared oxygen read returns, which measured coordinates
remain different, and what later outcome do they predict?}

Fix a run, a baseline sample, and a measurement declaration \(\mathcal M\).
The declaration assigns each quantity its assay, units, compartment,
organism or preparation, clock convention, and source reference. Let
\[
 R_0,A_0\in\pos,\qquad R,A\in\pos
\]
be the baseline and current readings of the declared oxidant and antioxidant
assays. Define
\begin{equation}\label{eq:normalized}
 r=\frac{R}{R_0},\qquad a=\frac{A}{A_0}.
\end{equation}
These are dimensionless, baseline-normalized assay indices. In particular,
an assay of glutathione concentration retains that identity when it occupies
the software field \code{scavenger}; the field name does not convert the assay
into a measurement of total scavenging capacity.

Write a numerical observation record as
\begin{equation}\label{eq:record}
 \mathcal R=(\mathcal M,\iota,n,t,O,R_0,A_0,R,A,\mathcal B,Y).
\end{equation}
Here \(\iota\) retains run, sample, and baseline identities; \(n\) is an
occurrence address; \(t\) is the visible clock value; \(O\) is the declared
oxygen read; \(\mathcal B\) retains entered bounds and other observations;
and \(Y\) is an optional, separately identified endpoint. One endpoint
definition belongs to one dataset. HIF isoforms, nuclear activity, serum EPO,
reticulocytes, and hemoglobin keep their individual measurement identities.

The implemented oxygen channel is one of nominal dry inspired partial
pressure, arterial partial pressure, or tissue partial pressure. For the
nominal dry-gas channel only, \(O=F_{\mathrm IO_2}P_{\mathrm{abs}}\).
The inspired fraction and absolute pressure remain separate inputs.
A return of this read is evaluated in its own compartment. The
oxygen-to-oxidant transfer is an additional empirical or modeling relation;
equation \eqref{eq:normalized} assigns none.

\section{A one-to-one coordinate interface}

\begin{theorem}[Positive assay-coordinate bijection]\label{thm:bijection}
The map
\[
 T:\pos^2\longrightarrow\pos^2,\qquad
 T(r,a)=(q,p)=\left(\frac r a,ra\right)
\]
is a bijection. Its inverse is
\begin{equation}\label{eq:inverse}
 r=\sqrt{pq},\qquad a=\sqrt{\frac p q},
\end{equation}
with the positive square roots. Consequently, for fixed positive \(R_0,A_0\),
\[
 (R,A)\longleftrightarrow(q,p)
 \quad\hbox{has inverse}\quad
 R=R_0\sqrt{pq},\qquad A=A_0\sqrt{p/q}.
\]
\end{theorem}
\begin{proof}
For a positive pair, \(pq=(ra)(r/a)=r^2\) and
\(p/q=(ra)/(r/a)=a^2\), so \eqref{eq:inverse} recovers the original
pair. Conversely, any \(p,q>0\) give positive values in
\eqref{eq:inverse}; their ratio is \(q\) and their product is \(p\).
Multiplication by the retained baselines recovers \(R,A\).
\end{proof}

The software name for \(p\) is \code{scale_squared}. Both coordinates
are required. For example, \((r,a)=(1,1)\) and \((2,2)\) have \(q=1\),
whereas their products are \(1\) and \(4\). An unchanged contrast admits
a changed magnitude.

\begin{corollary}[Reference-preserving record equivalence]\label{cor:record}
Replace the numerical pair \(R,A\) in \eqref{eq:record} by \(q,p\), leaving
every other component unchanged. This transformation is a bijection between
the admitted numerical records and their transformed records. Applying it
row by row, without reordering or discarding rows, is a bijection of finite
record sequences.
\end{corollary}
\begin{proof}
Use Theorem \ref{thm:bijection} independently at each row with its retained
baseline. The identity map on the remaining fields preserves their references
and order.
\end{proof}

This is the precise one-to-one claim. It concerns declared numerical records.
The original CSV and metadata bytes are retained separately, since numerical
parsing can identify different textual spellings of the same number.
The workbench augments observations with derived coordinates and records
input, code, and output hashes. A hash is an integrity reference to the
retained file; the file supplies the original observations.

\begin{proposition}[Units and changes of reference]\label{prop:units}
Multiplying \(R,R_0\) by any \(u_R>0\), and independently \(A,A_0\) by
\(u_A>0\), leaves \(r,a,q,p\) unchanged. If instead the baseline values
are replaced by \(R'_0=c_RR_0\), \(A'_0=c_AA_0\), then
\[
 r'=\frac r{c_R},\quad a'=\frac a{c_A},\quad
 q'=\frac{c_A}{c_R}q,\quad p'=\frac{p}{c_Rc_A}.
\]
\end{proposition}
\begin{proof}
Substitute in \eqref{eq:normalized} and \(T\).
\end{proof}
Multiplicative unit conversion preserves the read; a changed baseline is
a recorded change of reference. Assay offsets, background subtraction,
and censored readings require their declared measurement model.
The positive-domain interface does not insert an artificial positive value
for a zero or below-detection observation.

\section{Exact transport to the Metric Mirror operators}

Retain the root operators of Metric Mirror Symmetry \cite{metric}:
\begin{equation}\label{eq:root}
\begin{gathered}
 Q=\begin{pmatrix}1&1\\1&0\end{pmatrix},\qquad
 C_1=\begin{pmatrix}0&-1\\1&0\end{pmatrix},\qquad
 G=C_1Q-QC_1=\begin{pmatrix}-2&1\\1&2\end{pmatrix},\\
 H=\frac G{\sqrt5},\qquad S=\frac{2Q-\Id}{\sqrt5},\qquad
 \Phi=\frac{1+\sqrt5}{2}.
\end{gathered}
\end{equation}
They satisfy \(H^2=S^2=\Id\) and \(HS=-SH\).
Introduce the explicit assay adapter
\begin{equation}\label{eq:C}
 \Cr=\begin{pmatrix}\Phi&-1\\1&\Phi\end{pmatrix},
 \quad J=\begin{pmatrix}0&1\\1&0\end{pmatrix},
 \quad D=\operatorname{diag}(1,-1),
 \quad v=\Cr\begin{pmatrix}r\\a\end{pmatrix}.
\end{equation}
The subscript identifies the application coordinate map; it does not rename
the established orientation operator \(C_1\).

\begin{theorem}[Intertwining, mirror, and signed form]\label{thm:adapter}
The matrix \(\Cr\) is invertible, and
\begin{equation}\label{eq:conjugacies}
 \Cr^{-1}H\Cr=J,\qquad \Cr^{-1}S\Cr=D,\qquad
 \Cr^{T}G\Cr=\sqrt5(\Phi+2)J.
\end{equation}
For \(s\in\R\), let \(L_s=e^{sS}\). Then
\begin{equation}\label{eq:flow}
 L_s=\Cr\begin{pmatrix}e^s&0\\0&e^{-s}\end{pmatrix}\Cr^{-1},
 \quad HL_sH=L_{-s}=L_s^{-1},\quad L_s^TGL_s=G.
\end{equation}
On positive assay coordinates the mirror exchanges \(r,a\), and hence
maps \(q\) to \(q^{-1}\) while preserving \(p\).
\end{theorem}
\begin{proof}
Use \(\Phi^2=\Phi+1\). Then
\[
 \det\Cr=\Phi^2+1=\Phi+2>0,\qquad
 \Cr^{-1}=\frac1{\Phi+2}
   \begin{pmatrix}\Phi&1\\-1&\Phi\end{pmatrix}.
\]
Direct multiplication gives \(H\Cr=\Cr J\), \(S\Cr=\Cr D\),
and the last identity of \eqref{eq:conjugacies}.
Exponentiating \(S\Cr=\Cr D\) gives the first identity in
\eqref{eq:flow}. The equalities \(JDJ=-D\) and
\(\operatorname{diag}(e^s,e^{-s})^T J
\operatorname{diag}(e^s,e^{-s})=J\) give the other two.
Finally \(J(r,a)^T=(a,r)^T\).
\end{proof}

The form \(G\) has signature \((1,1)\), since \(G^2=5\Id\) and
\(\operatorname{tr}G=0\). On the admitted positive cone,
\begin{equation}\label{eq:productmetric}
 v^TGv=2\sqrt5(\Phi+2)ra.
\end{equation}
This is the pullback of a signed bilinear form. It assigns no physical
energy or concentration unit to that form.

\begin{theorem}[Magnitude times axial change]\label{thm:path}
For \(r,a>0\), define \(\zeta=\sqrt{ra}\) and
\(\eta=\tfrac12\log(r/a)\). Then
\begin{equation}\label{eq:representation}
 r=\zeta e^\eta,\quad a=\zeta e^{-\eta},\quad
 v=\zeta L_\eta\Cr\begin{pmatrix}1\\1\end{pmatrix}.
\end{equation}
For any two positive observations, with no dynamical assumption,
\begin{equation}\label{eq:path}
 v'=\frac{\zeta'}{\zeta}L_{\eta'-\eta}v.
\end{equation}
The mirror sends \(\eta\) to \(-\eta\), and \(L_s\) sends it to
\(\eta+s\) at fixed \(\zeta\).
\end{theorem}
\begin{proof}
The first two identities follow from the definitions and positivity.
Equation \eqref{eq:flow} gives \eqref{eq:representation}. Apply
\(L_{\eta'-\eta}L_\eta=L_{\eta'}\) and multiply by
\(\zeta'/\zeta\).
\end{proof}

The observations, or a separately declared model, supply
\(\eta'-\eta\) and \(\zeta'/\zeta\). The established golden step
\(L_{\log\Phi}=(Q+2\Id)/\sqrt5\) remains a coordinate step; it
supplies no biological rate. Equation \eqref{eq:path} explicitly retains
the magnitude change when \(ra\) changes.

\section{Exact sampled adaptation and non-erasing return}

Consider the declared adaptation model
\begin{equation}\label{eq:ode}
 \tau\frac{da}{dt}=r-a,\qquad \tau>0.
\end{equation}
Here \(r\) is a prescribed normalized oxidant drive. The equation is a
testable model of a chosen assay index. A measured oxygen channel and
a rule relating it to \(r\) retain their separate empirical assignments.

\begin{proposition}[Sampling and interval composition]\label{prop:sample}
If \(r\) is constant on an interval of length \(\Delta>0\), the exact
sampled solution is
\begin{equation}\label{eq:update}
 A_{\lambda,r}(a)=\lambda a+(1-\lambda)r,\qquad
 \lambda=e^{-\Delta/\tau}\in(0,1).
\end{equation}
It preserves positivity. For the same constant \(r\),
\begin{equation}\label{eq:composition}
 A_{\lambda_2,r}\circ A_{\lambda_1,r}
       =A_{\lambda_1\lambda_2,r}.
\end{equation}
\end{proposition}
\begin{proof}
Solving \eqref{eq:ode} gives
\(a(t+\Delta)=r+(a(t)-r)e^{-\Delta/\tau}\), which is
\eqref{eq:update}. It is a convex combination of positive values.
Composition gives
\(\lambda_2\lambda_1a+(1-\lambda_2\lambda_1)r\).
\end{proof}

One may choose a rational \(\lambda\in(0,1)\) directly for exact finite
arithmetic; its corresponding time constant is
\(\tau=-\Delta/\log\lambda\). Rational inputs and rational \(\lambda\)
give rational iterates. Combining \(m\) equal constant-input intervals
uses \(\lambda^m\). If the drive changes within the combined interval,
its segments must be retained. A change of time unit rescales
\(\Delta\) and \(\tau\) together and leaves \(\lambda\) unchanged.

\begin{theorem}[Pulse, jump, and finite recovery]\label{thm:return}
Let \(a=1\) initially. Apply the constant drive \(r=k>1\) for an
integer \(d\ge1\) intervals with \(0<\lambda<1\). Put removal of
that drive at time \(t=0\). Then
\begin{equation}\label{eq:pulse}
 a_d=1+(k-1)(1-\lambda^d),\qquad 1<a_d<k.
\end{equation}
At the jump, \(r(0^-)=k\), \(r(0^+)=1\), and
\(a(0^-)=a(0^+)=a_d\). After \(j\ge0\) recovery intervals,
\begin{equation}\label{eq:recovery}
 r_j=1,\qquad
 a_j=1+(a_d-1)\lambda^j,\qquad
 q_j=\frac1{1+(a_d-1)\lambda^j}<1.
\end{equation}
For every finite \(j\), the adaptation state differs from baseline.
As \(j\) increases, \(a_j\) decreases to one and \(q_j\) increases to one.
\end{theorem}
\begin{proof}
During the pulse, \(a_{n+1}-k=\lambda(a_n-k)\). Starting from
\(a_0=1\) gives \eqref{eq:pulse}; its inequalities follow from
\(0<\lambda^d<1\).
The solution of \eqref{eq:ode} is continuous at a bounded input jump,
so removal changes \(r\) and retains \(a_d\).
During recovery, \(a_{j+1}-1=\lambda(a_j-1)\); induction gives
\eqref{eq:recovery}. Its sign, monotonicity, and limit follow from
\(a_d-1>0\) and \(0<\lambda<1\).
\end{proof}

For \(k=2,\lambda=1/2,d=4\), the exact values are
\[
 (r,a,q)_{0^-}=(2,31/16,32/31),\quad
 (r,a,q)_{0^+}=(1,31/16,16/31),\quad
 (r,a,q)_{j=1}=(1,47/32,32/47).
\]
The two records at clock time zero have distinct occurrence addresses.
There is no elapsed adaptation interval between \(0^-\) and \(0^+\).

\begin{corollary}[The instantaneous drive does not determine the state]
\label{cor:nonfactor}
On the model class containing an unexposed baseline and the recovery
histories of Theorem \ref{thm:return}, there is no function \(f\) such
that \(a=f(r)\) at every observation.
\end{corollary}
\begin{proof}
The unexposed state has \(r=a=1\); a finite recovery state has \(r=1\)
and \(a>1\). The same argument of \(f\) would require two values.
\end{proof}
This proves failure of instantaneous-read factorization in the model.
The empirical question is whether specified measured records exhibit
the corresponding retained difference under their uncertainty bounds.

\section{History, uncertainty, and computational scope}

\begin{proposition}[A state and count do not recover an arbitrary history]
\label{prop:collision}
For \(\lambda=1/2\) and initial \(a=1\), the two drive histories
\((2,2)\) and \((3,3/2)\) both yield \(a=7/4\) after two intervals.
Returning each drive to one produces the same current pair \((r,a)\)
and the same interval count, while the histories are distinct.
\end{proposition}
\begin{proof}
The first history gives \(1\mapsto3/2\mapsto7/4\); the second gives
\(1\mapsto2\mapsto7/4\). Returning the drives changes both current
values of \(r\) to one without changing their adaptation state.
\end{proof}
Accordingly the occurrence record retains the raw sequence. The
read/state/occurrence distinction is inherited from Infinity One Zero
Lift \cite{iozl}; no eight-coordinate biological state is assumed.
The workbench creates a hashed snapshot. Retaining successive source
snapshots preserves the study record; the current coordinate pair
has the reconstruction scope of Theorem \ref{thm:bijection}.

\begin{proposition}[Positive interval enclosure]\label{prop:interval}
Suppose each of \(R,A,R_0,A_0\) belongs to its entered positive interval
\([X_-,X_+]\). Valid enclosing intervals are
\begin{align}
 r&\in[R_-/R_{0+},\,R_+/R_{0-}],&
 a&\in[A_-/A_{0+},\,A_+/A_{0-}],\label{eq:ra-bounds}\\
 q&\in\left[\frac{R_-A_{0-}}{A_+R_{0+}},
             \frac{R_+A_{0+}}{A_-R_{0-}}\right],&
 p&\in\left[\frac{R_-A_-}{R_{0+}A_{0+}},
             \frac{R_+A_+}{R_{0-}A_{0-}}\right].\label{eq:qp-bounds}
\end{align}
A baseline normalized by that identical baseline sample has the exact
self-ratios \(r=a=q=p=1\).
\end{proposition}
\begin{proof}
Each positive product or quotient increases in its numerator variables
and decreases in its denominator variables; the endpoints enclose the
rectangular domain. Additional dependence may narrow the attainable
set, so the enclosure remains valid. For an identical baseline sample,
numerator and denominator are the same variable and cancel.
\end{proof}

An interval excluding one resolves a change relative to the entered bounds.
An interval containing one leaves that sign unresolved. Equal point-valued
oxygen entries establish an exact return \emph{as entered}; overlapping
oxygen intervals establish compatibility with return. These are arithmetic
classifications. Probabilistic coverage requires its own measurement
specification. Exact rational arithmetic preserves entered numbers; it does
not remove assay uncertainty.

The mirror also has a precise dynamical boundary. The adaptation update
on pairs is \(\mathcal A_\lambda(r,a)=(r,\lambda a+(1-\lambda)r)\).
For \(\lambda=1/2\) at \((2,1)\),
\[
 J\mathcal A_{1/2}(2,1)=(3/2,2),
 \qquad \mathcal A_{1/2}J(2,1)=(1,3/2).
\]
Thus the coordinate mirror need not be a symmetry of the adaptation law.
Equation \eqref{eq:flow} retains the exact symmetry of its declared axial
family.

\section{The empirical adapter: primary fields and missing links}

An admitted empirical row needs a subject/run identity, its paired baseline,
an ordered sample time, the declared oxygen channel, both positive assay
readings, and their measurement descriptions. An endpoint is attached under
its own declaration. Additional covariates remain in the source record.
An available field maps to its named coordinate; an unmeasured field
remains unmeasured.

An endpoint acquired at a different time retains its own timed record,
linked through the subject/run and sampling protocol. It is not copied
into a redox row as if co-measured. The paired CSV interface requires
both assay readings; endpoint-only observations remain in their linked
source table until an explicit analysis defines the relevant cross-time
relation.

\begin{longtable}{@{}>{\raggedright\arraybackslash}p{0.18\linewidth}>{\raggedright\arraybackslash}p{0.38\linewidth}>{\raggedright\arraybackslash}p{0.39\linewidth}@{}}
\toprule
Primary reference & Reported observations & Interface assignment\\
\midrule
\endfirsthead
\toprule
Primary reference & Reported observations & Interface assignment\\
\midrule
\endhead
Donati et al.\ (2017) \cite{donati}
&
Serum HPF-reactive oxidant signal (RFU) and reported GSH
(\(\mu\mathrm M\)): baseline, end of two-hour exposure, and two hours
after return to individual baseline inspired fraction. Serum EPO:
baseline, 24 and 48 hours.
&
Same-sample oxidant and GSH readings can supply \code{ros} and
\code{scavenger} with their serum assay identities. Paired participant
values are required for \(q,p\); EPO is a separate later endpoint.\\[5pt]
Fratantonio et al.\ (2021) \cite{fratantonio}
&
PBMC transcription-factor and reduced-glutathione measurements after
a one-hour exposure; methods/figure times are 0.5, 3 and 24 hours.
ROS was not directly measured.
&
GSH can supply the declared antioxidant index and HIF-\(1\alpha\)
the specified endpoint. An observed paired ROS/GSH coordinate needs
an additional ROS measurement. Source compartment and timing
ambiguities stay attached to the import.\\[5pt]
Khalife et al.\ (2021) \cite{khalife}
&
Postoperative EPO, reticulocyte, and hemoglobin outcomes under repeated
oxygen variation and a comparison group with its own exposure history.
&
These supply separately identified outcomes and protocol context;
paired oxidant/antioxidant observations are not supplied.
Within-group and between-group contrasts retain separate estimands.\\
\bottomrule
\end{longtable}

Donati et al.\ report a transient oxidant rise and later GSH rise, giving
a direct empirical location for a retained-state comparison \cite{donati}.
The HPF signal is an assay-specific oxidant proxy; reported GSH retains
its concentration identity. Reduced-versus-total processing must be
resolved from the laboratory record. Fratantonio et al.\ connect exposure
to transcription-factor observations while leaving the paired oxidant
coordinate unmeasured \cite{fratantonio}. These are complementary interfaces,
not one merged participant trajectory.

The source ledger retains pressure-unit and assay-compartment qualifications,
and the 3-hour/4-hour wording discrepancy in Fratantonio et al.
A measured import must resolve these or carry explicit uncertainty.
The release contains no participant-level dataset from these studies.

The empirical connection to the matrices is exact at this declared interface:
\[
 \text{identified assay pair}
 \ \longleftrightarrow\ (r,a)
 \ \longleftrightarrow\ (q,p)
 \ \longleftrightarrow\ (\eta,\zeta)
 \ \longleftrightarrow\ v,
\]
with baselines and raw references retained. The arrows are invertible
coordinate maps on the positive domain. Oxygen-to-ROS transfer,
ROS/scavenger-to-PHD/HIF coupling, and downstream EPO response each
require their own observed relations.

Mechanism discrimination requires participant-level pairing, appropriate
oxygen and perfusion measurements, the declared assays, and time-resolved
PHD/HIF observations. Prespecified current-oxygen and history-dependent
models can be compared on held-out subjects or sequences. The HIF
observation law retains its measured sign and lag.
\textbf{Participant-level kinetic identification and ROS--HIF causation
are \OPEN.} Theorems \ref{thm:bijection}--\ref{thm:return} retain
their proved mathematical strength within their stated domains.

\section{Proof and executable-check crosswalk}

The proofs above establish the results on their declared domains.
The existing \code{verify_algebra.py} checks symbolic identities and exact
rational witnesses; \code{tests/test_audit.py} checks the observation adapter.
A fixture check has the finite scope stated below.

\small
\begin{longtable}{@{}>{\raggedright\arraybackslash}p{0.25\linewidth}>{\raggedright\arraybackslash}p{0.70\linewidth}@{}}
\toprule
Result & Existing checks and scope\\
\midrule
\endfirsthead
\toprule
Result & Existing checks and scope\\
\midrule
\endhead
Theorem \ref{thm:bijection}
&
\code{coordinate_reconstruction_r},
\code{coordinate_reconstruction_a}: symbolic positive-branch inverses.
\code{ratio_alone_loses_common_mode}: exact magnitude-loss witness.\\[3pt]
Proposition \ref{prop:units}
&
\code{ratio_unit_invariance},
\code{product_unit_invariance}: symbolic unit identities.
Baseline covariance is proved by substitution above.\\[3pt]
Theorem \ref{thm:adapter}
&
\code{adapter_invertible}, \code{adapter_S_diagonal},
\code{adapter_H_exchange}, \code{signed_metric_pullback},
\code{physical_state_pullback}: symbolic adapter identities.\\[3pt]
Theorems \ref{thm:adapter}, \ref{thm:path}
&
\code{mirror_reverses_axial_flow}, \code{axial_preserves_G},
\code{golden_step_diagonal}, \code{ratio_mirror_inversion},
\code{rational_boost_ratio}, \code{rational_boost_product}.
The full path identity is proved in Theorem \ref{thm:path}.\\[3pt]
Proposition \ref{prop:sample}; Theorem \ref{thm:return}
&
\code{pulse_closed_form_base}, \code{pulse_closed_form_recurrence},
\code{recovery_closed_form_base}, \code{recovery_closed_form_recurrence}:
symbolic recurrence identities.
\code{finite_timescale_rejects_zero_retention}: domain guard.\\[3pt]
Composition and jump
&
\code{constant_input_interval_composition},
\code{pulse_four_intervals_31_over_16},
\code{returned_ratio_16_over_31},
\code{first_recovery_47_over_32}: exact rational witnesses.
\code{finite_recovery_below_ratio_one}: supplied finite fixture;
the all-finite-time claim is proved in Theorem \ref{thm:return}.\\[3pt]
Proposition \ref{prop:collision}
&
\code{history_collision_retains_record_requirement}: the exact
two-history collision proved above.\\[3pt]
Proposition \ref{prop:interval}; record order
&
\code{test_baseline_self_ratios_cancel_even_with_uncertainty},
\code{test_uncertainty_overlap_remains_unresolved},
\code{test_interleaved_runs_and_tied_clock_keep_occurrence_order},
\code{test_assay_unit_conversion_preserves_all_normalized_quantities}:
implementation examples of the record contract.\\[3pt]
Dynamical boundary
&
\code{biochemical_update_is_not_assumed_mirror_symmetry}:
exact counterexample to imposing the coordinate exchange on adaptation.\\
\bottomrule
\end{longtable}
\normalsize

The receipts bind the executed algebra checker and its model dependency;
the data-audit receipt binds inputs, audit code, and output bytes.
These certify the construction checks and identify their implementation.
The empirical contract specifies how actual measurements enter the
construction while retaining their references and earned claim strength.

\begin{thebibliography}{9}
\footnotesize
\setlength{\itemsep}{1pt}
\setlength{\parskip}{0pt}
\setlength{\parsep}{0pt}
\bibitem{metric}
A.~V. Coppa. \emph{Metric Mirror Symmetry: Riemann's \(i\) Machine}.
Coppalini Architecture, 9 September 2026, Sections 1--7.
\href{https://doi.org/10.5281/zenodo.23478606}{doi:10.5281/zenodo.23478606}.

\bibitem{iozl}
A.~V. Coppa. \emph{Infinity One Zero Lift: The Normal Form of Projective Memory}.
Coppalini Architecture, working version 0.1, 11 November 2026.
Retained source manuscript, Sections 4, 7 and 8.

\bibitem{review}
M. Salvagno et al.
The Normobaric Oxygen Paradox---Hyperoxic Hypoxic Paradox:
A Novel Expedient Strategy in Hematopoiesis Clinical Issues.
\emph{International Journal of Molecular Sciences} 24(1), 82 (2023),
Section 3.2.
\href{https://doi.org/10.3390/ijms24010082}{doi:10.3390/ijms24010082}.

\bibitem{donati}
A. Donati et al.
Effects of short-term hyperoxia on erythropoietin levels and microcirculation
in critically ill patients: a prospective observational pilot study.
\emph{BMC Anesthesiology} 17, 49 (2017).
\href{https://doi.org/10.1186/s12871-017-0342-2}{doi:10.1186/s12871-017-0342-2}.

\bibitem{fratantonio}
D. Fratantonio et al.
Increasing Oxygen Partial Pressures Induce a Distinct Transcriptional
Response in Human PBMC: A Pilot Study on the Normobaric Oxygen Paradox.
\emph{International Journal of Molecular Sciences} 22(1), 458 (2021).
\href{https://doi.org/10.3390/ijms22010458}{doi:10.3390/ijms22010458}.

\bibitem{khalife}
M. Khalife et al.
Physiological and Clinical Impact of Repeated Inhaled Oxygen Variation
on Erythropoietin Levels in Patients After Surgery.
\emph{Frontiers in Physiology} 12, 744074 (2021).
\href{https://doi.org/10.3389/fphys.2021.744074}{doi:10.3389/fphys.2021.744074}.
\end{thebibliography}
\end{document}
